\documentclass[12pt,a4paper]{article}
\usepackage[T2A]{fontenc}
\usepackage[utf8]{inputenc}
\usepackage[russian]{babel}
\usepackage{amsmath,amssymb,geometry,hyperref}
\geometry{margin=18mm}
\hypersetup{hidelinks}
\setlength{\parindent}{0pt}
\setlength{\parskip}{5pt}
\begin{document}
\begin{center}\large\textbf{Лабораторная работа №3. Метод множителей Лагранжа}\end{center}
Борисюк Е. Г. Вариант 4.

\textbf{1. Постановка задачи}

Найти экстремумы функции $f(x,y)=xy$ при условии $x^2+4y^2=8$.
Ограничение записывается как $g(x,y)=x^2+4y^2-8=0$.

\textbf{2. Проверка нерегулярных точек}

При одном ограничении нерегулярная точка удовлетворяет $\nabla g=0$ и $g=0$:
\[\nabla g=(2x,8y)=(0,0)\quad\Longrightarrow\quad x=y=0.\]
Но $g(0,0)=-8\ne0$. Нерегулярных допустимых точек нет; метод Лагранжа
применим ко всем точкам условного экстремума.

\textbf{3. Функция Лагранжа и частные производные}

Знак функции Лагранжа выбирается как в конспектах практик:
\[L(x,y,\lambda)=xy+\lambda(x^2+4y^2-8).\]
\[L_x=y+2\lambda x,\qquad L_y=x+8\lambda y,\qquad L_\lambda=x^2+4y^2-8.\]

\textbf{4. Система необходимых условий}
\[\begin{cases}y+2\lambda x=0,\\x+8\lambda y=0,\\x^2+4y^2=8.\end{cases}\]
Если $x=0$, первое уравнение даёт $y=0$, что противоречит ограничению.
Аналогично $y=0$ невозможен. Поэтому $x\ne0$, $y\ne0$.
При $\lambda=0$ первые два уравнения также дают недопустимую пару $(0,0)$.

Из первого уравнения $y=-2\lambda x$. Подстановка во второе даёт
\[x-16\lambda^2x=x(1-16\lambda^2)=0.\]
Так как $x\ne0$, получаем
\[16\lambda^2-1=0,\qquad \lambda=-\frac14\quad\text{или}\quad\lambda=\frac14.\]
Других ветвей решения нет.
\newpage
\textbf{5. Решение системы и проверка кандидатов}

При $\lambda=-1/4$ имеем $y=x/2$. Тогда $2x^2=8$, откуда
$(x,y)=(2,1),(-2,-1)$.

При $\lambda=1/4$ имеем $y=-x/2$. Тогда снова $2x^2=8$, откуда
$(x,y)=(2,-1),(-2,1)$.

\begin{center}
\begin{tabular}{c|r|r|c|r|c}
$\lambda$ & $x$ & $y$ & $x^2+4y^2$ & $xy$ & Характер\\\hline
$-1/4$ & 2 & 1 & 8 & 2 & максимум\\
$-1/4$ & $-2$ & $-1$ & 8 & 2 & максимум\\
$1/4$ & 2 & $-1$ & 8 & $-2$ & минимум\\
$1/4$ & $-2$ & 1 & 8 & $-2$ & минимум
\end{tabular}
\end{center}
Во всех точках $x^2+4y^2=4+4=8$. Соответствующие множители также
обращают обе производные по $x,y$ в ноль.

\textbf{6. Характер экстремумов и полнота ответа}

Допустимое множество --- непустой замкнутый и ограниченный эллипс:
$|x|\le2\sqrt2$, $|y|\le\sqrt2$. Непрерывная функция достигает на нём
глобальных минимума и максимума. Все точки регулярны, поэтому каждый
экстремум входит в найденный полный список кандидатов.

Среди кандидатов получены только значения $2$ и $-2$. Следовательно,
глобальный максимум равен $2$, а минимум равен $-2$. Все четыре точки
являются и строгими локальными условными экстремумами: равенство
соответствующей глобальной границе достигается только в двух изолированных точках.

\textbf{Независимая проверка границ}
\[(x-2y)^2=x^2+4y^2-4xy=8-4xy\ge0,\]
\[(x+2y)^2=x^2+4y^2+4xy=8+4xy\ge0.\]
Значит, $-2\le xy\le2$ на всём эллипсе. Верхняя граница достигается при
$x=2y$, нижняя --- при $x=-2y$. Это подтверждает точки и значения.

\textbf{7. Ответ}
\[\boxed{f_{\max}=2\text{ в }(2,1),(-2,-1).}\]
\[\boxed{f_{\min}=-2\text{ в }(2,-1),(-2,1).}\]

\newpage
\textbf{8. Программная реализация}

В \texttt{main()} задаются исходные выражения $f=xy$ и
$g=x^2+4y^2-8$. Функция \texttt{solve\_lagrange} получает их и переменные.
Готовые координаты, производные и сведённое уравнение для множителя
в алгоритм не подставляются.

Программа вычисляет $\nabla g$ и решает систему $g=0$, $\nabla g=0$.
Затем составляет $L=f+\lambda g$, дифференцирует по $x,y,\lambda$,
решает построенную систему и подстановкой проверяет ограничение и
все три уравнения. Для варианта 4 нерегулярных точек нет.

Порядок метода и проверки реализованы самостоятельно. SymPy используется
для \texttt{diff}, \texttt{solve}, \texttt{simplify}; готовая функция
условной оптимизации не вызывается. Область реализации: две переменные,
одно равенство, изолированные действительные решения. Нерегулярные точки
и семейства решений требуют отдельного исследования.

\textbf{Характер точек по вторым производным}

В регулярной точке $t=(g_y,-g_x)^T$ --- ненулевой касательный вектор.
Программа строит Гессиан $L$ по исходным переменным и вычисляет $q=t^THt$:
\[H=\begin{pmatrix}2\lambda&1\\1&8\lambda\end{pmatrix},
\qquad t=\begin{pmatrix}8y\\-2x\end{pmatrix}.\]
При $q>0$ получается строгий локальный минимум, при $q<0$ --- строгий
локальный максимум. При $q=0$ второй порядок не даёт ответа.
Для $(2,1),(-2,-1)$ получено $q=-128$; для $(2,-1),(-2,1)$ --- $q=128$.

\textbf{9. Запуск и результат}
\begin{verbatim}
python -m pip install -r python/requirements.txt
python python/lagrange.py
\end{verbatim}
При запуске найдены четыре точки из ручного решения. Невязки равны нулю
в точной символьной арифметике. Наименьшее значение среди кандидатов $-2$,
наибольшее $2$. Глобальность для эллипса обоснована в разделе 6.
Код, выполненный Colab и вывод:
\url{https://github.com/Doney888/methods-of-optimization/tree/main/lab-03}.

\textbf{10. Рефлективный вывод}

Работа показывает связь условия параллельности градиентов с системой
Лагранжа. Автоматическое построение производных и проверка подстановкой
помогают контролировать ручные вычисления. Проверка вторых производных
определяет локальный характер, а доказательство существования экстремумов
и полноты кандидатов позволяет сделать глобальный вывод.

\end{document}
